\section*{Adopt a common FAR and DUPAC sample}
\textit{September 11, 2026. Agent-drafted record of Jacob's approved change.}

Jacob approved requiring usable values of both FAR and DUPAC for inclusion in
either density analysis. He clarified that each regression keeps its existing
distance restriction. Table 1 retains its citywide sample. The full working
paper now uses this common eligibility rule in the main figure, both placebo
figures, both donut figures, the boundary-type table, location-continuity table,
and stringency-score checks.

The analysis-data producer records the rule once as \texttt{density\_eligible}:
both density measures must be allowed, finite and positive. All 9,493 project
rows and their previously recorded values remain unchanged. The added field
marks 9,465 eligible projects within the 1,500-foot input population; regression
distance restrictions and complete controls determine each fitted sample.

The main sample has 4,023 projects overall and 863 multifamily projects for both
outcomes. Relative to the preceding version, this removes 11 projects from
DUPAC and leaves the main FAR sample unchanged. The estimated differences are
$-9.24\%$ for all-project FAR, $-13.70\%$ for all-project DUPAC, $-12.85\%$ for
multifamily FAR and $-17.36\%$ for multifamily DUPAC. The rounded multifamily
DUPAC effect changes from 18 to 17 percent in the abstract, results and
conclusion. Main differences remain significant at the 10 percent level or
better. Placebo estimates remain insignificant at 10 percent, and all donut
estimates remain negative. The location-continuity table still has no
significant differences at 10 percent. Table 1's data and formatted table are
unchanged, as are the existing permit and price results.

The own-project permit check was rebuilt on the common sample. Its coefficients
and standard errors match the main estimates at the table's reported precision.
Both score-gap restrictions retain their previous significance pattern. The
recorded project-permit matches and boundary-characteristic values are unchanged.

Jacob also approved explaining manual research directly. Appendix A now says
that construction years and lot or building areas sometimes appear inconsistent
in the Assessor files, that these records are supplemented with manual research
using permits, property records and other documented sources, and that the
replication files contain the decisions and supporting evidence. This replaces
the unverified 52-project total and single-spreadsheet description.

\smallskip
\noindent\textit{Reproduction:} The analysis-data Make target creates the common
eligibility field and standard report. The four density result tasks consume
that field. Run Make in \texttt{paper/} to build the complete working paper.
The analysis-data report has 9,493 unique project IDs and 66 columns; the only
added column is the common eligibility field. Unchanged Make dry runs for the
analysis and density tasks report no work. The completed paper and this entry
were compiled through Make and inspected as rendered PDFs.
